\section{The theorem and the structure of the proof}\label{sec:introduction}

How quickly does local resampling remove fluctuations in a densely coupled
spin system? For the Sherrington--Kirkpatrick model~\cite{SK1975}, the signs
of the interactions matter: each spin interacts weakly with every other spin,
but the sum of the absolute interaction strengths grows with the dimension.
We study heat-bath dynamics, which updates a spin from its conditional Gibbs
law, in the tradition of Glauber's local stochastic dynamics~\cite{Glauber1963}.
Our result is a dimension-independent Poincar\'e inequality throughout the
high-temperature phase, at each fixed inverse temperature below one. It controls every real observable on a single
event of disorder probability tending to one.

Fix $0<\beta<1$ and put $j=\beta^2$. Let $J$ be a real symmetric
$n\times n$ matrix with zero diagonal and independent off-diagonal entries
$J_{ij}\sim N(0,j/n)$, $i<j$. For a field $h\in\R^n$, write
\[
 \mu_h(x)=\frac{1}{Z_h}\exp\left\{\frac12x^{\mathsf T}Jx+h^{\mathsf T}x\right\},
 \qquad x\in\{-1,1\}^n.
\]
Let $P_i^h$ denote conditional expectation under $\mu_h$ given all spins
except the $i$th. Define the unscaled Dirichlet form
\begin{equation}\label{eq:dirichlet}
 \D_h(f)=\sum_{i=1}^n\E_{\mu_h}(f-P_i^h f)^2.
\end{equation}

We use a rate-one clock at each site in continuous time, giving the
generator $\mathcal L_h=\sum_i(P_i^h-I)$. A discrete attempt selects one
site uniformly and has transition operator $P_h=n^{-1}\sum_iP_i^h$.
Thus $\mathcal L_h=n(P_h-I)$, and the discrete gap is the continuous gap
divided by $n$. The gap of a reversible generator is the infimum of its
Dirichlet form divided by the variance, over all nonconstant functions.

\begin{theorem}\label{thm:main}
There is a finite constant $C_\beta$, depending only on $\beta$, such that
\[
 \Prob_J\left\{\Var_{\mu_0}(f)\le C_\beta\D_0(f)
       \text{ for every }f:\{-1,1\}^n\longrightarrow\R\right\}
 \longrightarrow1.
\]
Consequently, the discrete-time heat-bath chain that chooses a uniform site
has spectral gap at least $1/(C_\beta n)$ with probability tending to one.
\end{theorem}

The parameter $\beta$ is fixed before $n$ tends to infinity; the constant
$C_\beta$ may depend on that fixed choice. External fields enter the proof
through posterior Gibbs laws, while the theorem concerns $\mu_0$.
The diagonal convention is immaterial to the Gibbs law: because $x_i^2=1$,
adding any diagonal matrix to $J$ adds a constant to the Hamiltonian.
This observation will let us introduce a Gaussian diagonal for matrix
calculations and remove it again.

\subsection{Earlier results and methods}
The interval $\beta<1$ is the classical zero-field high-temperature
range. Aizenman, Lebowitz, and Ruelle~\cite[Proposition~2.1]{ALR1987}
proved that $n^{-1}\log Z_0$ converges in probability and in mean to
$\log2+\beta^2/4$, the annealed free-energy value, throughout this interval.
The dynamical question is whether local resampling also removes fluctuations at a rate independent
of the system size.

A basic route to such a bound controls conditional influences.
Wu~\cite{Wu2006} proved a Poincar\'e inequality under a
Dobrushin spectral-radius condition on a matrix of nonnegative
conditional influences.
In dense SK systems, stronger fixed-temperature estimates become possible
by exploiting the signed random interactions. Bauerschmidt and
Bodineau~\cite{BB2019} used a Gaussian field decomposition to prove a
high-temperature logarithmic Sobolev inequality; its full spin-flip energy
differs from the heat-bath form~\eqref{eq:dirichlet}.
Eldan, Koehler, and Zeitouni~\cite{EKZ2022} established a spectral
condition for the heat-bath gap, yielding the SK range $\beta<1/4$.
Adhikari, Brennecke, Xu, and Yau~\cite{ABXY2024} obtained constant
unscaled gaps for mixed $p$-spin models under a sufficiently small
weighted condition on their interaction coefficients.

For mixing in single-site discrete attempts, Anari, Jain, Koehler, Pham,
and Vuong~\cite{AJKPV2022} used entropic independence to prove an
$O(n\log n)$ bound in the SK range $\beta<1/4$.
Anari, Koehler, and Vuong~\cite{AKV2024} extended this order to an
inverse-temperature threshold approximately equal to $0.295$.
More recently, Wang~\cite[Theorem~1.1]{Wang2026} proved worst-start
mixing in $O_\beta(n\log(n/\varepsilon))$ attempts for every fixed
$\beta<1/2$, simultaneously over external fields. Boban, Li, and Oveis
Gharan~\cite[Theorem~1.5 and Corollary~1.6]{BLO2026} obtained a
constant unscaled gap and $O(n^2)$ zero-field mixing for
$\beta<1/2+\varepsilon_0$, with an absolute
$\varepsilon_0\ge5\cdot10^{-5}$. Theorem~\ref{thm:main} reaches
every fixed $\beta<1$ at zero field.

Equilibrium covariance estimates already cover this entire range.
El Alaoui and Gaitonde~\cite[Theorem~1.1]{EAG2024} proved that the
expected operator norm of the zero-field spin covariance matrix stays
bounded for every fixed $\beta<1$. Brennecke, Schertzer, Xu, and
Yau~\cite[Theorem~1.1]{BSXY2024} proved the operator-norm approximation
\[
 \left\|\Cov_{\mu_0}(x)-\bigl((1+\beta^2)I-J\bigr)^{-1}\right\|
 \longrightarrow0
 \qquad\text{in probability}.
\]
For $a\in\R^n$, the variance of the linear observable $a^{\mathsf T}x$
is $a^{\mathsf T}\Cov_{\mu_0}(x)a$. These covariance results therefore
control linear fluctuations. The all-functions heat-bath inequality
requires additional estimates; our proof obtains them at the random
posterior fields generated by Gaussian observations.

Our argument transfers a functional inequality along stochastic
localization, introduced by Eldan~\cite{Eldan2013}.
We use its Gaussian-observation representation developed by El Alaoui
and Montanari~\cite{ElAlaouiMontanari2021}. The variance-retention and
Dirichlet-form framework of Chen and Eldan~\cite{CE2025} explains how
posterior estimates can imply an inequality for the initial law.
Here an observation event of exponentially small probability can still
carry substantial variance for a particular test function. We control
its probability after weighting paths by the square of that function.
The entropy calculation is a stopped finite-mixture version of the
drift-energy representation; see Lehec~\cite{Lehec2013} and the work
of F\"ollmer discussed there.

Gaussian observations, planted laws, and approximate message passing
(AMP) are also combined in the SK sampling method of El Alaoui,
Montanari, and Sellke~\cite{EAMS2022}. Celentano~\cite{Celentano2024}
proved local convexity of the TAP free energy used by that sampler near
sufficiently late AMP iterates. Our argument needs invertibility at every
sufficiently accurate solution of the field equation along the observation
path. A scalar entropy inequality and a local-volume estimate give this
uniformity in Section~\ref{sec:stability}. The field recursions retain the
Onsager correction of the
Thouless--Anderson--Palmer equations~\cite{TAP1977} and are related to
Bolthausen's iterative construction~\cite{Bolthausen2014}. We initialize
them at a Gibbs spin sample and use exact conditional-spin identities
through a fixed finite depth, without requiring convergence of the
recursion. Finally, the terminal gap adapts Wang's signed two-spin
method~\cite{Wang2026}; Section~\ref{sec:terminal} proves the variant
with an entrywise-square cost and a quartic row error.

\subsection{Notation and elementary identities}
For a vector $v$, set $\av{v}=n^{-1}\sum_i v_i$ and let $D_v$ be the diagonal
matrix with diagonal $v$. Vector products and scalar functions of vectors are
understood coordinatewise. Unless another subscript is displayed, $\norm{\cdot}$
is Euclidean norm for vectors and operator norm for matrices;
$\norm{\cdot}_{\Frob}$ is Frobenius norm. A centered Gaussian orthogonal ensemble (GOE) matrix of scale $j$
has variances $j/n$ off diagonal and $2j/n$ on the diagonal.

If $x^{(i)}$ is obtained by flipping the $i$th spin, put
\[
 d_i f(x)=\frac{f(x)-f(x^{(i)})}{2x_i},\qquad
 (\nabla u)_{ij}=d_j u_i.
\]
Thus $d_i f$ is independent of $x_i$ and
\begin{equation}\label{eq:discrete-product}
 d_i(fg)=f\,d_i g+g\,d_i f-2x_i(d_i f)(d_i g).
\end{equation}
At a fixed field $h_0$ set
\[
 h^1=h_0+Jx,\qquad m^1=\tanh h^1,\qquad v^1=\sech^2 h^1.
\]
Since $J_{ii}=0$, $m_i^1$ and $v_i^1$ are respectively the conditional mean
and variance of $x_i$. In particular,
\begin{equation}\label{eq:dirichlet-gradient}
 \D_{h_0}(f)=\E_{\mu_{h_0}}\sum_i v_i^1(d_i f)^2,\qquad
 \norm{G}_*^2=\E_{\mu_{h_0}}G^2+\D_{h_0}(G).
\end{equation}
For probability laws $\nu,\mu$, their relative entropy is
$D(\nu\Vert\mu)=\int\log(\dd\nu/\dd\mu)\,\dd\nu$ when $\nu\ll\mu$,
and $+\infty$ otherwise.

Constants are independent of $n$; dependence on a fixed temperature,
construction depth, or positive stability margin is recorded at its use.
Their values may change between displays. The prerequisites are Gaussian
concentration and comparison, Gaussian integration by parts, and stochastic
calculus. We prove the model-specific estimates in the sections below.

\subsection{The proof mechanism}
There are two parts to the argument: prove useful inequalities for
posterior Gibbs laws, then transfer them to the unobserved law. A long but
fixed observation makes most individual spins predictable. The corresponding
small conditional variances, combined with a two-spin curvature inequality,
give a terminal gap. To recover the initial variance, we need a directional
covariance estimate along the path and a mean estimate under square
reweighting. The former bounds the loss of variance; the latter prevents
that variance from accumulating on exceptional paths.

The posterior estimates come from finite Onsager-corrected recursions.
Exact conditional-spin identities control their residuals. A telescoping
identity for the squared increments then selects two consecutive small
increments without requiring convergence of the recursion. Uniform matrix
estimates handle
the resulting adaptive diagonal coefficients under an explicit inverse
margin. To compare these approximate fields with an exact root, stability
must hold throughout the small-residual region and persist when the
diagonal variance matrix is replaced by the nearby response coefficients
used in the comparison. A scalar entropy estimate controls a Gaussian
integral, and
a local-volume lower bound turns that integral control into exclusion of
unstable approximate roots. All tolerances and finite depths are chosen
before the large-$n$ limit.

\subsection{The three inputs}
We observe a spin sample in independent Gaussian noise:
\begin{equation}\label{eq:observation-process}
 X\sim\mu_0,\qquad Y(t)=tX+B(t),\quad 0\le t\le T,
\end{equation}
where $B$ is standard $n$-dimensional Brownian motion independent of $X$.
Conditional on the observations up to $t$, the spin law is $\mu_{Y(t)}$.
The proof uses the following three inputs, established in the later sections.

\begin{enumerate}
\item \textbf{Stability along ordinary observation paths.}
Proposition~\ref{prop:stable-fields} gives a deterministic notion of a good
field such that, for typical $J$, all fields $Y(t)$ are good except on an event
of observation probability at most $e^{-a n}$, where $a>0$ is fixed.

\item \textbf{Two estimates at a good field.}
For every $G$, Proposition~\ref{prop:directional} proves
\begin{equation}\label{eq:input-directional}
 \norm{\Cov_{\mu_h}(G,x)}^2
 \le C\bigl(\Var_{\mu_h}(G)+\D_h(G)\bigr).
\end{equation}
For $N=\E_{\mu_h}f^2>0$ and $\nu=f^2\mu_h/N$,
Proposition~\ref{prop:weighted-mean} proves, for each sufficiently small fixed $\delta>0$,
\begin{equation}\label{eq:input-weighted}
 \frac{\norm{\E_\nu x-t_*(h)}}{\sqrt n}
 \le C\delta+\frac{C_\delta}{\sqrt n}
                 \left(1+\frac{\D_h(f)}{N}\right).
\end{equation}
Here $t_*(h)=\tanh r(h)$, where $r(h)$ is the unique solution of
\[
 r-h+\{j(1-n^{-1}\|\tanh r\|^2)I-J\}\tanh r=0.
\]
Existence and uniqueness at the fields under consideration are part of
Lemma~\ref{lem:root}. The constant $C$ is independent of $\delta$; $C_\delta$ may depend
on it. Taking $f=1$ gives the same comparison for $\E_{\mu_h}x$.

\item \textbf{A terminal gap.}
For a sufficiently large fixed $T$, Proposition~\ref{prop:terminal-gap}
gives $\Var_{\mu_{Y(T)}}f\le2\D_{Y(T)}(f)$ outside an event of ordinary
observation probability at most $e^{-a n}$.
\end{enumerate}

The transfer uses the two estimates for different purposes. Directional
covariance controls posterior variance loss, while square-reweighted means
bound the entropy of the changed observation law. Their combination controls
the contribution of the bad paths before the terminal gap is applied.
Section~\ref{sec:observation} proves this implication first, so the later
technical sections have a precise target.

\subsection{Organization}
After the observation transfer in Section~\ref{sec:observation},
Section~\ref{sec:matrices} constructs the matrix diagnostics used by the
adaptive field calculations. Sections~\ref{sec:entropy}
and~\ref{sec:stability} identify the scalar equality law and establish
stability along observations. Section~\ref{sec:residuals} develops the
finite recursions and their weak residual rules; Section~\ref{sec:consequences}
uses them to prove the two posterior estimates. Section~\ref{sec:terminal}
closes the argument with the fixed-time terminal gap.
